% Compile with XeLaTeX. Content converted from kontrolna-stem-2-varianty-bez-vidpovidey.docx.
\documentclass[11pt,a4paper]{article}
\usepackage{fontspec}
\usepackage[ukrainian]{babel}
\usepackage[margin=18mm]{geometry}
\usepackage{array,tabularx}
\setmainfont{Arial}
\setlength{\parindent}{0pt}
\setlength{\parskip}{6pt}
\renewcommand{\arraystretch}{1.65}
\begin{document}
\begin{center}\Large\textbf{Діагностична робота з курсу STEM}\par\large І семестр, 5 клас\end{center}

На основі навчального плану модулів «Я у Всесвіті», «Я так бачу!», «Під знаком STEM»\par

Прізвище, ім'я \rule{0.65\linewidth}{0.4pt}\par

Клас \rule{0.25\linewidth}{0.4pt}    Дата \rule{0.25\linewidth}{0.4pt}\par

\section*{Варіант 1}

\subsection*{1. Тестові завдання}

\begin{minipage}{\linewidth}1. Який прилад використовують для спостереження за небесними тілами?\par\begin{tabularx}{\linewidth}{@{}XX@{}}А. Телескоп & Б. Мікроскоп\\ В. Фотоальбом & Г. Лупа\end{tabularx}\end{minipage}\par

\begin{minipage}{\linewidth}2. Що таке орбіта?\par\begin{tabularx}{\linewidth}{@{}XX@{}}А. Космічний костюм & Б. Шлях руху небесного тіла\\ В. Малюнок планети & Г. Назва ракети\end{tabularx}\end{minipage}\par

\begin{minipage}{\linewidth}3. Який об'єкт належить до космічної техніки?\par\begin{tabularx}{\linewidth}{@{}XX@{}}А. Герб & Б. Піктограма\\ В. Ракета & Г. Печатка\end{tabularx}\end{minipage}\par

\begin{minipage}{\linewidth}4. Що допомагає створювати рухомі зображення?\par\begin{tabularx}{\linewidth}{@{}XX@{}}А. Анімація & Б. Орбіта\\ В. Телескоп & Г. Скафандр\end{tabularx}\end{minipage}\par

\begin{minipage}{\linewidth}5. Який із наведених прикладів є штучним знаком?\par\begin{tabularx}{\linewidth}{@{}XX@{}}А. Веселка & Б. Темні хмари\\ В. Дорожній знак & Г. Слід лапи\end{tabularx}\end{minipage}\par

\begin{minipage}{\linewidth}6. Для чого космонавту потрібен скафандр?\par\begin{tabularx}{\linewidth}{@{}XX@{}}А. Для прикрашання & Б. Для захисту в космосі\\ В. Для малювання & Г. Для подачі сигналів\end{tabularx}\end{minipage}\par

\subsection*{2. Теоретичне питання}

Поясни, яку роль відіграють кольори у графіці, символах і повсякденному житті людини.\par

\rule{0.96\linewidth}{0.4pt}\par

\rule{0.96\linewidth}{0.4pt}\par

\rule{0.96\linewidth}{0.4pt}\par

\rule{0.96\linewidth}{0.4pt}\par

\rule{0.96\linewidth}{0.4pt}\par

\newpage\textbf{Варіант 1 — продовження}\par

\subsection*{3. Практичне завдання 1}

Заповни таблицю про космос.\par

\begin{tabularx}{\linewidth}{|>{\raggedright\arraybackslash}X|>{\raggedright\arraybackslash}X|>{\raggedright\arraybackslash}X|}\hline

\textbf{Об'єкт / поняття} & \textbf{Що це або для чого використовується} & \textbf{До якої теми належить} \\ \hline

Телескоп &  &  \\ \hline

Орбіта &  &  \\ \hline

Ракета &  &  \\ \hline

Скафандр &  &  \\ \hline

Планета &  &  \\ \hline

Космічна станція &  &  \\ \hline

\end{tabularx}\par

\subsection*{4. Практичне завдання 2}

Склади план розробки моделі ракети. Запиши 5-7 послідовних кроків.\par

1. \rule{0.65\linewidth}{0.4pt}\par

2. \rule{0.65\linewidth}{0.4pt}\par

3. \rule{0.65\linewidth}{0.4pt}\par

4. \rule{0.65\linewidth}{0.4pt}\par

5. \rule{0.65\linewidth}{0.4pt}\par

6. \rule{0.65\linewidth}{0.4pt}\par

7. \rule{0.65\linewidth}{0.4pt}\par

\newpage

Прізвище, ім'я \rule{0.65\linewidth}{0.4pt}\par

Клас \rule{0.25\linewidth}{0.4pt}    Дата \rule{0.25\linewidth}{0.4pt}\par

\section*{Варіант 2}

\subsection*{1. Тестові завдання}

\begin{minipage}{\linewidth}1. Який приклад є природним знаком?\par\begin{tabularx}{\linewidth}{@{}XX@{}}А. Смайлик & Б. Буква\\ В. Темні хмари & Г. Печатка\end{tabularx}\end{minipage}\par

\begin{minipage}{\linewidth}2. Що належить до графічних знакових систем?\par\begin{tabularx}{\linewidth}{@{}XX@{}}А. Герб & Б. Планета\\ В. Орбіта & Г. Скафандр\end{tabularx}\end{minipage}\par

\begin{minipage}{\linewidth}3. Що таке піктограма?\par\begin{tabularx}{\linewidth}{@{}XX@{}}А. Небесне тіло & Б. Умовний малюнок-знак\\ В. Космічний апарат & Г. Вид фарби\end{tabularx}\end{minipage}\par

\begin{minipage}{\linewidth}4. Яке з наведеного пов'язане з фотографією?\par\begin{tabularx}{\linewidth}{@{}XX@{}}А. Фотоапарат & Б. Телескоп\\ В. Герб & Г. Печатка\end{tabularx}\end{minipage}\par

\begin{minipage}{\linewidth}5. Що допомагає передавати інформацію в символічній формі?\par\begin{tabularx}{\linewidth}{@{}XX@{}}А. Знаки і символи & Б. Лише планети\\ В. Лише кольори & Г. Тільки ракети\end{tabularx}\end{minipage}\par

\begin{minipage}{\linewidth}6. До якого модуля належить тема фотографії та анімації?\par\begin{tabularx}{\linewidth}{@{}XX@{}}А. Людина - природа & Б. Людина - образ\\ В. Людина - знак & Г. Людина - техніка\end{tabularx}\end{minipage}\par

\subsection*{2. Теоретичне питання}

Поясни, як кольори допомагають людині передавати настрій, значення та інформацію в малюнках, символах і фотографіях.\par

\rule{0.96\linewidth}{0.4pt}\par

\rule{0.96\linewidth}{0.4pt}\par

\rule{0.96\linewidth}{0.4pt}\par

\rule{0.96\linewidth}{0.4pt}\par

\rule{0.96\linewidth}{0.4pt}\par

\newpage\textbf{Варіант 2 — продовження}\par

\subsection*{3. Практичне завдання 1}

Заповни таблицю про знаки і символи.\par

\begin{tabularx}{\linewidth}{|>{\raggedright\arraybackslash}X|>{\raggedright\arraybackslash}X|>{\raggedright\arraybackslash}X|}\hline

\textbf{Знак / символ} & \textbf{Що означає або для чого використовується} & \textbf{До якої теми належить} \\ \hline

Герб &  &  \\ \hline

Печатка &  &  \\ \hline

Піктограма &  &  \\ \hline

Дорожній знак &  &  \\ \hline

Буква &  &  \\ \hline

Емблема &  &  \\ \hline

\end{tabularx}\par

\subsection*{4. Практичне завдання 2}

Склади план розробки моделі фотоапарата. Запиши 5-7 послідовних кроків.\par

1. \rule{0.65\linewidth}{0.4pt}\par

2. \rule{0.65\linewidth}{0.4pt}\par

3. \rule{0.65\linewidth}{0.4pt}\par

4. \rule{0.65\linewidth}{0.4pt}\par

5. \rule{0.65\linewidth}{0.4pt}\par

6. \rule{0.65\linewidth}{0.4pt}\par

7. \rule{0.65\linewidth}{0.4pt}\par
\end{document}
